\section{Data and Model Selection}
\subsection{Dataset and audit}
The experiments use the public TrialGPT criterion-annotation dataset \cite{jin2024trialgpt}. The audit retained \DatasetAnnotationCount{} annotations from \DatasetPatientCount{} patients and \DatasetTrialCount{} trials after dropping empty criterion text. GPT-4-generated explanations and labels were excluded at load time; training targets use only expert eligibility labels and expert evidence sentence IDs. The fixed split is patient-disjoint, with seed 42 and counts shown in Table~\ref{tab:split}.

\begin{table}[t]
\centering
\caption{Patient-disjoint dataset partition sizes.}
\label{tab:split}
\begin{tabular}{lrrrr}
\toprule
Split & Patients & Annotations & Inclusion & Exclusion \\
\midrule
TRAIN & 37 & 602 & 215 & 387 \\
VAL & 8 & 206 & 66 & 140 \\
TEST & 8 & 206 & 75 & 131 \\
\bottomrule
\end{tabular}
\end{table}

\begin{table}[t]
\centering
\caption{LoRA training configuration read from the versioned YAML files.}
\label{tab:hyperparameters}
\begin{tabular}{lrrrrrr}
\toprule
Model & Max tokens & Epochs & Batch x accum. & Rank/alpha & Learning rate & Seed \\
\midrule
1.5B & 640 & 3 & 4 x 4 & 16/32 & 2.0e-04 & 42 \\
0.5B & 640 & 3 & 4 x 4 & 16/32 & 2.0e-04 & 42 \\
\bottomrule
\end{tabular}
\end{table}


The split is frozen by a manifest containing the dataset revision and patient assignments. All prompt, training, and model comparisons use the same split; the held-out test partition is accessed only by the tagged final evaluation.

\subsection{Models and supervised fine-tuning}
Qwen2.5-1.5B-Instruct is the primary target and Qwen2.5-0.5B-Instruct is the size comparison and same-tokenizer draft \cite{yang2025qwen25}. LoRA updates query, key, value, output, and MLP projections while the base weights remain frozen \cite{hu2022lora}. The training target is a compact JSON object with evidence IDs followed by a criterion-type-specific label. The patient note is placed before the criterion; explanations are generated from cited text by a deterministic template, not learned from GPT-4 columns.

Training histories record optimization loss and validation loss by epoch. They do not contain generation-based validation F1 for every checkpoint. Therefore, the report compares the completed final adapters on the already-run fixed validation evaluations and does not claim an epoch was selected by generation F1 (D-006).
The versioned training settings are summarized in Table~\ref{tab:hyperparameters}.
